\documentclass{beamer}
\usepackage{graphicx}
\newcommand{\sframe}[2]{\begin{frame}\frametitle{#1}\vspace{-1cm}~#2\end{frame}}
\newcommand{\sitem}{\\ - }
\newcommand{\ssitem}{\\\hphantom{8mm}- }
\newcommand{\sep}{\\~}

\begin{document}

\sframe{idnai-esp32 package: Motivation}{
\sitem Fast and easy development of connected objects.
\sep
\sitem Allows projects:
\ssitem cheap (ESP32 is a all in one cheap chip);
\ssitem easy (only specific handle functions to define);
\ssitem fast (elementary C programming).
\sep
\sitem The EPS32:
\centerline{\includegraphics[height=0.3\textheight]{esp32-view.jpg}}
}

\sframe{The package functionalities}{
\sitem Implements a generic setup and loop mechanism with:
\ssitem WIFI automatic setup;
\ssitem Web HTML page and Web service simple managements;
\ssitem USB console management.
\sep
\sitem Makes available:
\ssitem A fill-in web-service with from user handler;
\ssitem Basic I/O functions documentation;
\ssitem Time measurement functions;
\ssitem Periodic or delayed functions.
\sep
\sitem Helps development with:
\ssitem Elaborated echo functions for debug;
\ssitem User I/O web interface and I/O web service;
\ssitem Documented examples.
}

\sframe{How to ?}{
Everything is documented on the home page.
\sep
\sitem Once:
\ssitem Install the Arduino IDE on your desktop computer.
\ssitem Install your ESP32 project, from `idnai-esp32`.
\sep
\sitem Then:
\ssitem Develop your hardware once or step by step. 
\ssitem Connect the board on the desktop via USB.
\ssitem Starts the compilation/load/check/correct work loop.
}

\end{document}
